\section{Reproducibility checklist}
\label{app:reproducibility}

\subsection{Implementation components}

The implementation steps corresponding to the results reported in the paper
are:

\begin{itemize}
    \item \path{1_build_legal_causal_graph.py}: construct the legal graph;
    \item \path{2_build_causal_memory.py}: construct rule/event memory and
    the FAISS index;
    \item \path{3_multi_hop_causal_retriever.py}: retrieve seeds, expand, and
    rank causal paths;
    \item \path{causal_core/schema.py} and
    \path{causal_core/legal_scm.py}: implement canonical rules and structural
    inference;
    \item \path{4_counterfactual_verification.py}: perform query-aware
    verification, structural intervention, and path ablation;
    \item \path{5_generate_final_answer.py}: select evidence and generate
    answers;
    \item \path{5_5_generate_pipeline_predictions.py}: provide the batch
    runner for Steps~3--5;
    \item \path{7_compute_evaluation_metrics.py}: compute metrics and
    group-level diagnostics.
\end{itemize}

\subsection{Versions and principal settings}

\begin{table}[h]
    \centering
    \small
    \caption{Fixed versions and settings for the main experiment.}
    \label{tab:repro-settings}
    \begin{tabular}{@{}ll@{}}
        \toprule
        \textbf{Component} & \textbf{Setting} \\
        \midrule
        Benchmark & \texttt{2.0}, $250$ samples \\
        Step 4 & \texttt{4.1-query-aware-legal-scm} \\
        Step 5 & \texttt{5.1-query-aware-grounded-answer} \\
        Batch runner & \texttt{2.1-query-aware-answer} \\
        Evaluator & \texttt{2.0-current-pipeline-schema} \\
        Embedding & \texttt{BAAI/bge-m3} \\
        Retrieval hop / CF hop & $2$ / $3$ \\
        Answer provider & \texttt{extractive} \\
        Structural intervention & $\doop(M=\stateF)$ \\
        \bottomrule
    \end{tabular}
\end{table}

The complete thresholds are reported in
Table~\ref{tab:retrieval-hyperparameters},
Table~\ref{tab:verification-hyperparameters}, and
Table~\ref{tab:generation-hyperparameters}.

\subsection{Inputs and paired artifacts}

The principal inputs are:

\begin{itemize}
    \item \path{data/blhs_multihop_benchmark_250.json};
    \item \path{data/blhs_rules_final_all_normalized.json};
    \item the graph, causal memory, and FAISS index shared by both modes.
\end{itemize}

The prediction and evaluation directories used to report the results are:

\begin{itemize}
    \item \path{data/pipeline_predictions_structural_scm_v41_step5_v51.json};
    \item \path{data/pipeline_predictions_path_ablation_v41_step5_v51.json};
    \item \path{evaluation_results/structural_scm_v41_step5_v51/};
    \item \path{evaluation_results/path_ablation_v41_step5_v51/}.
\end{itemize}

The older artifacts in
\path{evaluation_metrics_path_ablation_v41/} and the predictions generated
by runner \texttt{2.0-primary-path-aware} are not part of the main paired
experiment and were not used to compute the tables in the paper.

\subsection{Sample batch command}

With \texttt{MODE} set in turn to \texttt{structural\_scm} and
\texttt{path\_ablation}, the batch was run using the following template:

\begin{verbatim}
python 5_5_generate_pipeline_predictions.py ^
  --benchmark data/blhs_multihop_benchmark_250.json ^
  --provider extractive ^
  --counterfactual-mode MODE ^
  --rules data/blhs_rules_final_all_normalized.json ^
  --output data/pipeline_predictions_MODE_v41_step5_v51.json ^
  --jsonl-output data/pipeline_predictions_MODE_v41_step5_v51.jsonl ^
  --errors-output data/pipeline_errors_MODE_v41_step5_v51.json ^
  --run-log data/pipeline_run_log_MODE_v41_step5_v51.json ^
  --work-dir data/pipeline_intermediate_MODE_v41_step5_v51 ^
  --checkpoint-every 1 --resume
\end{verbatim}

On Linux, the \texttt{\^{}} line-continuation character is replaced with a
backslash. Step 7 was run with the same benchmark and corresponding
predictions, with \texttt{--k-values 1 3 5 10} and \texttt{--strict}.

\subsection{Minimum trace schema}

Each prediction must retain at least the following provenance fields:

\begin{itemize}
    \item sample ID, query, and benchmark version;
    \item retrieval result, candidate path, and primary path ID;
    \item query analysis, claim type, and answer intent;
    \item final verification decision and verification score;
    \item factual/intervention state when structural SCM is used;
    \item evidence ID, rule ID, event ID, and article citation;
    \item Step 4, Step 5, runner, and evaluator versions;
    \item processing time and errors, if any.
\end{itemize}

\subsection{Pre-publication checks}

The following items must still be added from the exact server and repository
that produced the artifacts; they must not be inferred:

\begin{itemize}
    \item operating system, Python version, CPU, RAM, and GPU;
    \item exact embedding-model revision and dependency lockfile;
    \item immutable Git commit and public repository URL;
    \item hashes of the benchmark, rule corpus, graph, memory, and predictions;
    \item random seed and cache state;
    \item licenses for the code, data, and model;
    \item corresponding-author information, funding, and venue.
\end{itemize}

\todo{Fill in the hardware/server information, commit, artifact hash, and
author email before submitting the manuscript.}
